% !TEX root = ../main.tex
\chapter{Composite Operations}
\label{ch:composite_ops}

Not every operation in \Lucid's surface is a primitive C++ kernel.
A substantial set --- including most of the elementwise convenience
functions, the higher-level reductions, and the
quality-of-life utilities --- is implemented in pure Python on top
of the primitive surface. These are the \emph{composite
operations}, and they live under \code{lucid/\_ops/composite/}.

The composite layer is the design decision that keeps the C++
engine focused on operations whose performance materially matters
while delegating convenience operations to a layer that is easier
to extend and audit. This chapter describes what lives in the
composite layer, what the rules for composition are, and why some
operations migrate from composite to primitive over time.

\secref{sec:composite_ops:why} motivates the composite layer.
\secref{sec:composite_ops:rules} describes the rules.
\secref{sec:composite_ops:examples} gives examples.
\secref{sec:composite_ops:numpy_free} treats the numpy-free
discipline.
\secref{sec:composite_ops:promotion} discusses when a composite
should be promoted to a primitive.

\section{Why a Composite Layer}
\label{sec:composite_ops:why}

A composite operation expresses one ML operation as a sequence of
primitive ones. \code{softmax(x, dim)} can be expressed as
\code{exp(x - max(x, dim, keepdim=True))} divided by the
keepdim-true sum of that exponential. \code{logsumexp(x, dim)} can
be expressed as \code{max(x, dim, keepdim=True) + log(sum(exp(x -
max(x, dim, keepdim=True))))}. \code{cross\_entropy(logits, target)}
can be expressed as the negative log of the softmax of the logits,
indexed by the target.

Each of these compositions could in principle be a C++ kernel ---
and indeed, for the most performance-critical cases
(\code{logsumexp}, \code{softmax}, \code{cross\_entropy}), \Lucid\
has C++ primitives that match the composite's semantics. But the
composite expression has three properties that make it worth
keeping around even when a primitive exists:

\begin{enumerate}
  \item \textbf{Implementation reference.} The composite is the
    correctness reference against which the primitive is parity-
    tested. A divergence between the two is a bug in the
    primitive, the composite, or both.
  \item \textbf{Backward derivation.} The composite's backward is
    constructed automatically by autograd from its constituent
    primitives' backwards. This is essentially free; deriving the
    primitive's backward by hand is not.
  \item \textbf{Extensibility.} Adding a new composite op is a few
    lines of Python in a single file; adding a primitive is the
    nine-file pattern of \chref{ch:op_registry}.
\end{enumerate}

The trade-off is performance: a composite of \(n\) primitives
dispatches \(n\) times through the framework, each dispatch
costing the per-op overhead of
\chref{ch:backend_dispatch}~Section~\ref{sec:backend_dispatch:perf}.
For ops that are called millions of times in a training step (the
inner-loop softmaxes of an attention layer, for example), the
overhead can matter and we promote to primitive.

\section{Rules for the Composite Layer}
\label{sec:composite_ops:rules}

The composite layer follows three rules:

\begin{enumerate}
  \item \textbf{Numpy-free.} A composite uses only \Lucid\
    primitives and the Python standard library. No
    \code{import numpy}, no \code{import scipy}. The Hard Rule H4
    (\chref{ch:bridge_boundaries}) is enforced here.
  \item \textbf{Pure functions of tensors.} A composite takes
    tensor inputs and returns tensor outputs (plus an optional
    Python-typed metadata like a list of ints for shapes). It does
    not touch global state, does not mutate inputs, does not have
    side effects.
  \item \textbf{Autograd-clean.} A composite's autograd graph is
    constructed by the primitives it calls. The composite itself
    does not register a custom \code{Function} subclass.
\end{enumerate}

The third rule has an exception for numerically delicate composites
that need a custom backward to avoid the
divide-by-tiny-number issue that the natural composite would
have. \code{logsumexp} is the canonical example: its forward
expression is naturally numerically stable, but the autograd-
constructed backward through \code{log} would expose the
divide-by-sum that the forward avoided. The composite's actual
implementation overrides the backward with a closed-form softmax,
which is the analytic gradient.

\section{Examples}
\label{sec:composite_ops:examples}

The composite layer contains roughly 60 operations. We give a few
representative examples.

\subsection{softmax}
\label{ssec:composite_ops:softmax}

\code{softmax(x, dim)} is the canonical example:

\begin{lstlisting}[style=lucid,language=LucidPython]
def softmax(x, dim=-1):
    m = x.max(dim=dim, keepdims=True).values
    e = (x - m).exp()
    s = e.sum(dim=dim, keepdims=True)
    return e / s
\end{lstlisting}

Six lines, four primitive calls (\code{max}, \code{sub},
\code{exp}, \code{sum}, \code{div}). The autograd graph is
constructed by the primitives; the backward is automatic. The
numerical stability comes from the \code{x - m} subtraction inside
the \code{exp}, which is the standard trick.

A C++ primitive \code{softmax} also exists, with the same
numerics. The dispatcher prefers the primitive for the GPU path
when the input is large enough to amortise the dispatch overhead;
the composite remains as the correctness reference.

\subsection{cross\_entropy}
\label{ssec:composite_ops:cross_entropy}

\code{cross\_entropy(logits, target, reduction='mean')}:

\begin{lstlisting}[style=lucid,language=LucidPython]
def cross_entropy(logits, target, reduction='mean'):
    log_probs = lucid.log_softmax(logits, dim=-1)
    nll = -log_probs.gather(dim=-1, index=target.unsqueeze(-1)).squeeze(-1)
    if reduction == 'mean':
        return nll.mean()
    if reduction == 'sum':
        return nll.sum()
    return nll
\end{lstlisting}

A higher-level composite that itself uses another composite
(\code{log\_softmax}). The autograd flows through the entire chain
correctly.

\subsection{l2\_normalize}
\label{ssec:composite_ops:l2_normalize}

\code{l2\_normalize(x, dim)}:

\begin{lstlisting}[style=lucid,language=LucidPython]
def l2_normalize(x, dim=-1, eps=1e-12):
    norm = (x * x).sum(dim=dim, keepdims=True).sqrt().clamp(min=eps)
    return x / norm
\end{lstlisting}

The \code{eps}-clamp prevents division by zero for zero-norm input
rows.

\subsection{logsumexp with custom backward}
\label{ssec:composite_ops:logsumexp}

\code{logsumexp(x, dim)}:

\begin{lstlisting}[style=lucid,language=LucidPython]
class _LogSumExp(lucid.autograd.Function):
    @staticmethod
    def forward(ctx, x, dim):
        m = x.max(dim=dim, keepdims=True).values
        out = (m + (x - m).exp().sum(dim=dim, keepdims=True).log())
        ctx.save_for_backward(x, out)
        ctx.dim = dim
        return out.squeeze(dim=dim)

    @staticmethod
    def backward(ctx, grad_out):
        x, out = ctx.saved_tensors
        # Gradient is softmax: exp(x - logsumexp(x))
        return (x - out).exp() * grad_out.unsqueeze(ctx.dim), None

def logsumexp(x, dim=-1):
    return _LogSumExp.apply(x, dim)
\end{lstlisting}

This is the exception to the autograd-clean rule. The custom
backward returns the closed-form softmax rather than letting
autograd derive it through \code{log}, which would expose the
divide-by-sum that the forward avoided.

\section{The Numpy-Free Discipline}
\label{sec:composite_ops:numpy_free}

The composite layer must not import numpy. The reasoning is the
same as for the rest of \code{lucid/}
(\chref{ch:bridge_boundaries}): numpy imports inject an external
dependency into the numerical path, potentially with their own
dtype or rounding conventions.

\subsection{What replaces numpy in the composite layer}
\label{ssec:composite_ops:replace_numpy}

The patterns:

\begin{itemize}
  \item \code{np.zeros(...)} $\to$ \code{lucid.zeros(...)}.
  \item \code{np.broadcast\_to(...)} $\to$
    \code{lucid.broadcast\_to(...)}.
  \item \code{np.array(...)} $\to$ \code{lucid.tensor(...)}.
  \item Mathematical constants (\code{np.pi}, \code{np.e}) $\to$
    Python's \code{math.pi}, \code{math.e}, or hardcoded floats.
  \item Bitwise integer manipulation $\to$ standard Python
    operators.
\end{itemize}

The standard library covers most needs. The cases that needed
careful replacement during the 3.0.2 numpy-independence pass were
documented in \code{retro-3-0-2-numpy-free-standalone}.

\subsection{The discipline at code review}
\label{ssec:composite_ops:review_discipline}

The static check that enforces H4 runs on every PR and rejects any
\code{import numpy} (or scipy, etc.) in a composite file. A
contributor who needs a missing utility in the composite layer is
expected to either reach for a Python standard-library equivalent
or to add a primitive that the composite can then use.

\section{When to Promote a Composite to a Primitive}
\label{sec:composite_ops:promotion}

The composite layer is intended to be the first place a new op is
implemented. Promotion to primitive is reserved for cases that
justify the extra implementation cost.

\subsection{Performance criterion}
\label{ssec:composite_ops:perf_criterion}

The performance criterion: if the composite's per-call cost on a
representative training workload exceeds 10\,\% of the total step
time, the op is a candidate for promotion. This is rare but
happens for the very-inner-loop ops:

\begin{itemize}
  \item \code{softmax}, \code{log\_softmax}: promoted because
    every Transformer attention layer calls them many times per
    step.
  \item \code{layer\_norm}, \code{group\_norm}: promoted because
    they appear once per layer and the per-call overhead matters.
  \item \code{cross\_entropy}: promoted because losses are
    typically called once per batch and the small per-call
    overhead would otherwise dominate small-batch training.
\end{itemize}

For ops that are called less frequently (a handful of
\code{l2\_normalize} calls per step), the composite is fine.

\subsection{The promotion procedure}
\label{ssec:composite_ops:promotion_procedure}

When promotion is decided:

\begin{enumerate}
  \item Implement the C++ primitive following
    \chref{ch:op_registry}'s nine-file pattern.
  \item Wire the dispatcher to prefer the primitive for the
    relevant device.
  \item Keep the composite as a fallback (for the cases the
    primitive cannot handle, e.g., unusual dtypes the primitive
    does not implement) and as the parity reference.
  \item Add an explicit parity test that the primitive and the
    composite produce bit-equivalent outputs for representative
    inputs.
\end{enumerate}

The composite is not removed. Its continued existence in the test
suite is what protects against subtle divergence between the two
implementations as either evolves.

\subsection{When the composite stays}
\label{ssec:composite_ops:composite_stays}

Most composites stay forever:

\begin{itemize}
  \item Pure convenience functions (\code{l2\_normalize},
    \code{stable\_logit}, \code{soft\_threshold}) that wrap a few
    primitives.
  \item Algorithm-specific utilities (\code{moving\_average} for
    optimisers, \code{exponential\_moving\_variance} for AMP
    scaler).
  \item Reference implementations of operations that an
    application-level user might want to override
    (\code{label\_smoothed\_cross\_entropy} for some training
    regimes).
\end{itemize}

The composite layer is the engineering counterpart to the
primitive layer: the latter for hot paths, the former for
breadth.

\section{Catalogue of Composition Patterns}
\label{sec:composite_ops:catalogue}

Beyond the few examples already shown, the composite layer
exhibits a small set of recurring patterns. We enumerate them with
a representative composite for each.

\subsection{Numerical-stability shift}
\label{ssec:composite_ops:numerical_shift}

The pattern: subtract a max (or min) before an exponential to
prevent overflow.

Examples: \code{softmax}, \code{log\_softmax}, \code{logsumexp},
the inner reduction of \code{cross\_entropy}.

The pattern always has the form
\code{exp(x - x.max(dim=d, keepdims=True))}, with the corresponding
backward derived analytically (typically softmax-shaped).

\subsection{Broadcast-then-reduce}
\label{ssec:composite_ops:broadcast_reduce}

The pattern: lift two tensors into a higher-dim space via broadcast,
operate elementwise, then reduce back.

Examples: pairwise distance computations
(\code{(x[:, None, :] - y[None, :, :]).pow(2).sum(dim=-1)}),
contraction-free attention scores
(\code{(q.unsqueeze(-2) * k.unsqueeze(-3)).sum(dim=-1)} ---
though this is suboptimal; use \code{matmul}).

The pattern is convenient but memory-expensive: the intermediate
tensor has shape $N \cdot M$, which can dwarf the inputs.

\subsection{Reduce-then-normalise}
\label{ssec:composite_ops:reduce_norm}

The pattern: compute a reduction, broadcast its result back, and
divide.

Examples: \code{softmax} (subtract max, exp, normalise by sum),
\code{l2\_normalize} (sqrt(sum-of-squares), divide), \code{mean
centring} (sum, divide by N, subtract).

The composite-then-primitive trajectory has happened to several
members of this pattern: \code{softmax}, \code{layer\_norm}, and
\code{l2\_normalize} all started as composites and have C++
primitives now for the GPU hot path.

\subsection{Iterated apply}
\label{ssec:composite_ops:iterated}

The pattern: apply the same op repeatedly with a varying parameter.

Examples: \code{moving\_average} (one apply per timestep),
\code{exponential\_moving\_variance} (parameter is a decay rate),
the inner loops of certain second-order optimisers
(\code{lbfgs}'s line search).

These compositions are nontrivial to lift to a C++ primitive
because the iteration depth is data-dependent and the
per-iteration parameters vary. They remain in Python; the
per-iteration overhead is amortised by the work inside each
iteration.

\subsection{Algorithm-specific decomposition}
\label{ssec:composite_ops:algo_specific}

The pattern: a published algorithm with a specific decomposition
that does not generalise.

Examples: \code{label\_smoothed\_cross\_entropy} (one-hot smooth +
KL divergence), \code{focal\_loss} (cross-entropy times a
focal-weighting), \code{huber\_loss} (piecewise quadratic /
linear).

These are decomposition-specific and the user often wants to read
the algorithm; the Python composite is more pedagogically valuable
than a C++ primitive would be.

\section{When the Composite Is the Whole Story}
\label{sec:composite_ops:whole_story}

Most composites stay as composites forever, because their per-call
cost is not on the critical path of any workload. The full list of
\emph{permanent} composites at the time of writing:

\begin{itemize}
  \item Convenience wrappers: \code{l2\_normalize},
    \code{l1\_normalize}, \code{stable\_logit}, \code{soft\_threshold},
    \code{cosine\_similarity}.
  \item Loss helpers: \code{focal\_loss}, \code{huber\_loss},
    \code{label\_smoothed\_cross\_entropy}, \code{dice\_loss},
    \code{ctc\_loss\_padded}.
  \item Statistical utilities: \code{moving\_average},
    \code{exponential\_moving\_variance},
    \code{exponential\_moving\_mean}.
  \item Algorithm-specific: \code{adam\_update},
    \code{lamb\_update}, \code{lbfgs\_line\_search} (used internally
    by optimisers).
\end{itemize}

These ops collectively account for less than 1\,\% of training
step time on any workload we have measured; the composite path is
correct and the primitive path is not worth building.

\section{Composite Layer Hygiene}
\label{sec:composite_ops:hygiene}

A short list of disciplines that have kept the composite layer
maintainable across the 3.x evolution.

\subsection{One file per family}
\label{ssec:composite_ops:one_file}

Each composite-op family lives in a single file under
\code{lucid/\_ops/composite/}: activations in \code{activation.py},
elementwise in \code{elementwise.py}, reductions in
\code{reductions.py}, and so on. The file boundary corresponds
roughly to the C++ category split, which makes the parity
relationship easier to maintain.

\subsection{Public API at the right level}
\label{ssec:composite_ops:api_level}

A composite is exposed at the level the user is expected to call
it. \code{softmax} appears at \code{lucid.softmax} (Tier 1, free
function), at \code{lucid.nn.functional.softmax} (Tier 2,
functional surface), and as the building block of
\code{lucid.nn.Softmax} (Tier 2, stateful module). Less-used
composites like \code{l2\_normalize} appear only at
\code{lucid.nn.functional}.

The Tier 1 surface is curated to remain compact; not every
composite warrants top-level exposure.

\subsection{Parity-test discipline}
\label{ssec:composite_ops:parity_discipline}

Every composite has a parity test against the reference framework
under the path
\code{lucid/test/parity/ops/composite/}\allowbreak. The test
exercises forward, backward (via \code{gradcheck}), and AMP
behaviour. A
composite that lacks a parity test is rejected at review (per the
parity-test policy in \chref{ch:testing}).

The discipline matters because composites can drift: a primitive
that the composite calls may change behaviour (e.g., a schema
version bump), and without a parity check the drift goes
unnoticed.

\section{Summary and Forward References}
\label{sec:composite_ops:summary}

The composite layer in \code{lucid/\_ops/composite/} expresses
roughly 60 operations in pure Python on top of \Lucid\ primitives.
The layer is numpy-free, autograd-clean (with the documented
exception for numerical-stability custom backwards), and the first
place a new op should be implemented. Promotion to primitive is
reserved for the hot-path operations whose per-call overhead
matters; the composite version stays as the parity reference.

This chapter closes \partref{part:ops_kernels}. The remaining
parts of the book cover increasingly higher-level concerns:
\partref{part:autograd} treats the differentiation engine that
underlies the autograd machinery used throughout the op layer,
\partref{part:math_subsystems} treats the mathematical
subpackages, \partref{part:neural_networks} treats the neural
network library, \partref{part:graph_compile} treats the compile
path, and so on through model zoo, performance engineering, and
production engineering.

\begin{lucidnote}
For the user: the composite layer is the layer the user is most
likely to extend. Writing a custom op that combines a few
primitives is one Python file; writing a primitive is a multi-day
C++ project. For most extensions, the composite path is the
right choice.
\end{lucidnote}
